\section{The optimization problem and the main result}
The coherence of a family of unit vectors is the largest modulus of an
inner product between distinct vectors. Minimizing coherence is a
packing problem for complex lines: it asks how well a prescribed number
of lines can be separated in a given dimension.
All inner products are conjugate-linear in the first variable:
$\langle v,w\rangle=v^*w$. For integers $1\leq n<N$ define
\[
 \mu_{\mathbb C}(n,N)=
 \min_{\substack{v_1,\ldots,v_N\in\mathbb C^n\\\|v_i\|=1}}
 \max_{i\ne j}|v_i^*v_j|.
\]
The minimum exists: the product of $N$ unit spheres is compact and the
maximum of the finitely many continuous inner-product moduli is continuous.

We study the case $n=N-2$. The associated Gram matrices have a kernel
of dimension at least two; a two-dimensional subspace of that kernel
provides the geometric representation used below. When $4\mid N$, the
classical attaining construction corresponds to equal replication of
four equiangular lines in this representation~\cite{BC,Glazyrin}.
For the remaining cardinalities, equal replication is impossible.
We construct a candidate using balanced integer multiplicities and
deform the four-line configuration accordingly. We prove that this
candidate is globally optimal among all vector families for every
$N\ge180$, and determine all equality cases.

The kernel argument also applies to matrices that need not initially
be Gram matrices. We therefore consider the two matrix problems
\[
 a_N=\min_{\substack{A\in\mathbb C^{N\times N}\\A_{ii}=1,\ \rank A\leq N-2}}
       \max_{i\ne j}|A_{ij}|,\qquad
 b_N=\inf_{\substack{A\in\mathbb C^{N\times N}\\A_{ii}=1,\ \rank A=N-2}}
       \max_{i\ne j}|A_{ij}|.
\]
For $N\geq3$ the first minimum also exists. The all-ones matrix is feasible
with objective one, so we can restrict to the compact set with all
off-diagonal moduli at most one. Diagonal and rank-at-most conditions are
closed. Moreover $a_N>0$, because a diagonal-one matrix with zero
off-diagonal entries has rank $N$. No compactness of the exact-rank set is
assumed; attainment of $b_N$ will follow from the construction.

Set $\alpha=(1+\sqrt3)/4$. A four-tuple
$\boldsymbol m=(m_1,m_2,m_3,m_4)$ of positive integers with sum $N$ is \emph{balanced} if its
entries differ by at most one. Such a tuple is unique up to permutation.
Put $p_a=m_a/N$ and denote by $\lambda_{\boldsymbol p}$ the root specified
in the following theorem.

\begin{theorem}\label{thm:main}
For every integer $N\geq180$ and every balanced four-tuple
$\boldsymbol m$ of positive integers with sum $N$, the equation
\begin{equation}\label{eq:main-root}
 2\lambda=\sum_{a=1}^4\frac{p_a(1-3p_a)}{\lambda-2p_a}
\end{equation}
has a unique root in $(2/3,\alpha]$, and
\begin{equation}\label{eq:main-value}
 \mu_{\mathbb C}(N-2,N)=a_N=b_N=
 \frac1{N\lambda_{\boldsymbol p}-1}.
\end{equation}
The exact-rank infimum is attained. Every minimizing matrix in the
rank-at-most problem is, up to index permutation, diagonal unitary
switching and entrywise conjugation, the balanced four-block Gram matrix
in Proposition~\ref{prop:construction}. In particular it is positive semidefinite,
has rank $N-2$, and all off-diagonal entries have the same modulus.
Vector realizations differ additionally by a common unitary isometry.
\end{theorem}

For $4\mid N$ the value is $1/(\alpha N-1)$, the known divisible case.
For $N\equiv2\pmod4$ the formula becomes
\begin{equation}\label{eq:residue-two}
 \mu_{\mathbb C}(N-2,N)=
 \frac4{N+\sqrt{3N^2-8}-4}.
\end{equation}

\begin{corollary}\label{cor:asymptotic}
For each fixed $r\in\{0,1,2,3\}$, as $N\longrightarrow\infty$ through
$N\equiv r\pmod4$,
\[
 \mu_{\mathbb C}(N-2,N)=\frac1{\alpha N-1}
 +\frac{r(4-r)}{4\sqrt3\,\alpha^2}\,N^{-3}+O(N^{-4}).
\]
For $r\ne0$ the correction is positive.
\end{corollary}
\begin{proof}
Write $p_a=1/4+\epsilon_a$, $\sum_a\epsilon_a=0$ and
$\sum_a\epsilon_a^2=r(4-r)/(4N^2)$. Define
$\Phi_{\boldsymbol p}(\lambda)=4[2\lambda-
\sum_ap_a(1-3p_a)/(\lambda-2p_a)]$ and put $h=\alpha-1/2$.
Polynomial division gives
\[
 \Phi_{\boldsymbol p}(\alpha)
 =4\sum_a\frac{\epsilon_a^2}{h-2\epsilon_a},\qquad
 \Phi'_{(1/4,\ldots,1/4)}(\alpha)=\frac{4\sqrt3}{h}.
\]
The derivative is positive near uniform weights; Taylor's formula gives
$\alpha-\lambda_{\boldsymbol p}
=\sum_a\epsilon_a^2/\sqrt3+O(N^{-3})$.
Substitution in \eqref{eq:main-value} proves the assertion.
\end{proof}

The correction in Corollary~\ref{cor:asymptotic} measures the cost of the
integer multiplicities. It vanishes when equal replication is possible
and is positive in each of the other residue classes.

The proof begins with a trace inequality relating the coherence of a
matrix to an energy associated with a two-dimensional subspace of its
left kernel. A matrix that matches or improves on the construction must
have kernel energy close to the classical upper bound. We use the
resulting shared deficit bounds to place every kernel
column in one of four clusters and to restrict their integer counts.
A local comparison then controls both motion of the cluster centers and
variation within each cluster, including the complex phases. Its strict
inequality forces an extremizer to coincide with the corresponding
four-block candidate. Integer balancing selects the optimal counts,
and equality in the trace inequality recovers the full matrix.

The explicit construction and matrix argument are given in
Sections~\ref{sec:construction} and~\ref{sec:proof}. The quantitative
localization and local comparison use rational inequalities on closed
domains, whose certificates are described in Section~\ref{app:certificates}.
The threshold $180$ is sufficient; its minimality remains open.
